%! Rational Functions and Zeros — Unit 1, Lesson 1.8 — Slides (Beamer)
\documentclass[aspectratio=169, 11pt]{beamer}
\usepackage{precalculus-beamer}   % provides \navyheader{} and \sectionlabel[color]{}

\begin{document}

% TITLE SLIDE
{
\setbeamercolor{background canvas}{bg=navy}
\begin{frame}[plain]
  \vspace{2.8cm}\hspace{0.55cm}%
  \begin{minipage}{0.9\textwidth}
    {\color{goldacc}\bfseries\small LESSON 1.8}\\[0.5cm]
    {\color{white}\bfseries\LARGE Rational Functions and Zeros}\\[1.2cm]
    {\color{skymid}\small Precalculus~$\cdot$~Unit 1: Polynomial and Rational Functions}
  \end{minipage}
\end{frame}
}

% HOOK SLIDE
\begin{frame}[t, plain]
  \navyheader{Hook: Profit Per Unit}
  \sectionlabel[goldacc]{Think About It}

  A company models its \textbf{profit per unit} as
  \[
    P(x) = \frac{x^2 - 9}{x - 1}, \qquad x > 0,\; x \ne 1
  \]
  where $x$ is thousands of units produced.

  \bigskip
  \begin{block}{Discussion}
    \begin{itemize}\setlength{\itemsep}{6pt}
      \item At what production level(s) is the profit per unit \textbf{exactly zero}?
      \item For which values of $x$ is the profit \textbf{positive}?
      \item Why can't we evaluate $P(1)$?
    \end{itemize}
  \end{block}

  \bigskip
  \textit{Today: we find where rational functions equal zero and where they are positive/negative.}
\end{frame}

% ZEROS OF RATIONAL FUNCTIONS
\begin{frame}[t, plain]
  \navyheader{Zeros of Rational Functions}

  \begin{block}{Key Principle (EK 1.8.A.1)}
    For $r(x) = \dfrac{p(x)}{q(x)}$, \; $r(x) = 0$ when:
    \begin{enumerate}\setlength{\itemsep}{4pt}
      \item $p(x) = 0$ \quad (numerator equals zero), AND
      \item $q(x) \ne 0$ \quad ($x$ is in the domain)
    \end{enumerate}
  \end{block}

  \bigskip
  \textbf{Example:} Find the zeros of $\displaystyle r(x) = \dfrac{x^2 - 4}{x - 3}$.

  \medskip
  \begin{itemize}\setlength{\itemsep}{5pt}
    \item Domain: $x \ne 3$ (denominator zero)
    \item Set numerator $= 0$: $x^2 - 4 = 0 \Rightarrow (x-2)(x+2) = 0$
    \item Solutions $x = 2$ and $x = -2$: both satisfy $x \ne 3$, so both are in the domain
    \item \textbf{Zeros: $x = 2$ and $x = -2$}
    \item Note: $x = 3$ is \textbf{NOT} a zero --- it is undefined there
  \end{itemize}
\end{frame}

% SIGN CHART METHOD
\begin{frame}[t, plain]
  \navyheader{Sign Chart Method for Rational Inequalities}

  \textbf{Solve:} $\displaystyle\frac{x - 1}{x + 2} \geq 0$

  \medskip
  \begin{block}{4 Steps}
    \begin{enumerate}\setlength{\itemsep}{4pt}
      \item \textbf{Find critical values:} \quad Numerator zero: $x = 1$. \; Denominator zero: $x = -2$.
      \item \textbf{Mark on number line:} \quad open circle at $x = -2$; closed circle at $x = 1$
      \item \textbf{Test one value per interval:}
            \begin{itemize}
              \item $x = -3$: $(-)/(-)= + \Rightarrow (+)$
              \item $x = 0$: $(-)/( +) \Rightarrow (-)$
              \item $x = 2$: $(+)/(+) \Rightarrow (+)$
            \end{itemize}
      \item \textbf{Include intervals where sign is $(+)$; include zero $x=1$; exclude $x=-2$}
    \end{enumerate}
  \end{block}

  \medskip
  \textbf{Solution:} $(-\infty,\,-2)\cup[1,\,\infty)$
\end{frame}

% CLASS PRACTICE
\begin{frame}[t, plain]
  \navyheader{Class Practice}
  \sectionlabel[redacc]{Try It}

  Work with your partner on the following two problems.

  \bigskip
  \textbf{Problem 1.} Find all zeros of $\displaystyle f(x) = \dfrac{x^2 - 5x + 6}{x + 1}$.
  \begin{itemize}\setlength{\itemsep}{4pt}
    \item What is the domain restriction?
    \item Factor the numerator and find numerator zeros.
    \item Which numerator zeros are in the domain?
  \end{itemize}

  \bigskip
  \textbf{Problem 2.} Solve $\displaystyle\frac{x + 3}{x - 2} \leq 0$ using a sign chart.
  \begin{itemize}\setlength{\itemsep}{4pt}
    \item List critical values and mark the number line.
    \item Test one value per interval.
    \item Write the solution in interval notation.
  \end{itemize}
\end{frame}

\end{document}
